\section{The actual partial-strip quotient at polynomial weight zero}
\label{sec:hardy-partial-strip}

This section computes the actual Hilbert closure and its convolution
trace when the polynomial weight parameter is exactly $s=0$.
The strip need not contain every completed zero.  No uniform distance
of its interior zeros from its edges is assumed.  The argument uses
the canonical generator of MT-M18, a Hardy-space closure calculation,
and actual compressed multiplication operators.  Its conclusion is
strictly about the partial quotient; it supplies no return map to the
full smooth quotient when an exterior block is lost, as MA-HM-003
already proves.

\subsection{The norm and the actual closed ideal}

Fix the primitive unitary Hecke function, its entire pole-removed
completion $\Xi=\Xi_\eta$, and
\[
 \mathcal I_\eta=\overline{\psi_\eta*\mathcal E}^{\,\mathcal E},
 \qquad \mathcal M\psi_\eta=\Xi_\eta
\]
of MT-M18.  Let $a\in\mathbb R$, $u<v$, $L=v-u$, and $c=(u+v)/2$.
Throughout this section the polynomial weight parameter is $s=0$.
Write
\[
 \begin{aligned}
 H_{\max}&=L^2(\mathbb R,\max(e^{2ux},e^{2vx})\,dx),\\
 H_\Sigma&=L^2(\mathbb R,(e^{2ux}+e^{2vx})\,dx),\\
 J&=\overline{T_a\mathcal I_\eta}^{\,H_{\max}}
   =\overline{T_a\mathcal I_\eta}^{\,H_\Sigma},
       &T_af(x)&=e^{ax}f(x),\\
 Z_{\max}&=H_{\max}/J,\qquad Z_\Sigma=H_\Sigma/J,\qquad
 \Pi=\{z:u<\operatorname{Re}z<v\}.
 \end{aligned}
 \tag{MA-HY.1}
\]
The same set $J$ occurs because
\[
       \|f\|_{\max}\le\|f\|_\Sigma\le\sqrt2\,\|f\|_{\max}.
 \tag{MA-HY.2}
\]
The identity on representatives consequently induces a bounded
invertible map between the two actual quotients.  It intertwines
translation and convolution.  In particular it preserves trace
class and trace by bounded similarity; it need not identify the
two orthogonal projections onto the respective complements of $J$.

Let $\mathscr H^2(\Pi)$ denote the translation-invariant Hardy space
with norm
\[
 \|F\|_{\mathscr H^2}^2
 =\sup_{0<\delta<L/2}\frac1{2\pi}\int_{\mathbb R}
       \bigl(|F(u+\delta+it)|^2+
             |F(v-\delta+it)|^2\bigr)\,dt .
 \tag{MA-HY.3}
\]
This is the unweighted boundary-norm Hardy space, rather than the
space obtained by composition alone with a disk conformal map.
The Mellin transform is a unitary map
\[
             \mathcal M:H_\Sigma\longrightarrow\mathscr H^2(\Pi).
 \tag{MA-HY.4}
\]
Indeed, Plancherel on each vertical line identifies the integral
in (MA-HY.3) with
\[
 \int_{\mathbb R}|f(x)|^2
       \bigl(e^{2(u+\delta)x}+e^{2(v-\delta)x}\bigr)\,dx .
\]
Its integrand increases as $\delta$ decreases to zero, since the
parenthesis equals $2e^{2cx}\cosh((L-2\delta)x)$.
Monotone convergence proves the norm identity.  Conversely, if
$F$ satisfies (MA-HY.3), take inverse Fourier transforms $f_\sigma$
of $t\mapsto F(\sigma+it)$ on the interior lines.  The
Cauchy--Riemann equation gives, in distributions,
$\partial_\sigma f_\sigma=x f_\sigma$.  Testing first on compact
$x$ intervals and solving this scalar equation gives
$f_\sigma=e^{\sigma x}f$ for one distribution $f$, which is a
locally square-integrable function because any one $f_\sigma$
is in $L^2$.  The uniform bound and Fatou's lemma as
$\sigma\to u,v$ give $f\in H_\Sigma$.  The interior Mellin
integral is absolutely convergent by Cauchy--Schwarz against
$(e^{2ux}+e^{2vx})^{-1}$, and equals $F$ by Fourier uniqueness.
This proves surjectivity, and also the $L^2$ boundary limits.

For later use choose the explicit conformal maps
\[
 q(w)=c+\frac{2L}{\pi}\arctan w,\qquad
 \varphi(z)=\tan\frac{\pi(z-c)}{2L},\qquad
       q'(w)=\frac{2L}{\pi(1+w^2)} .
 \tag{MA-HY.5}
\]
Thus $q:\mathbb D\to\Pi$, and the two infinite ends correspond to
$w=i$ and $w=-i$.  With any analytic square root of $q'$,
\[
     \mathcal U F(w)=q'(w)^{1/2}F(q(w))
        :\mathscr H^2(\Pi)\longrightarrow H^2(\mathbb D)
 \tag{MA-HY.6}
\]
is unitary.  The normalization follows from the boundary change
of variables $|q'(e^{i\theta})|\,d\theta=dt$, with both edges
included.  The onto statement, including the Hardy-space
condition at the two ends, is the rotated and rescaled
$p=2$ case of Bakan--Kaijser (2007), Theorems~2.1--2.2 and
Corollary~2.1, pp.~349--350, with proofs in \S4, pp.~352--355.
In particular the half-derivative factor in (MA-HY.6) is required.

There is also a direct verification of the precise onto statement,
including its end condition.  The reproducing kernel of the
Mellin image in (MA-HY.4) is
\[
\begin{aligned}
 K_\Pi(z,\zeta)
 &=\int_{\mathbb R}
     \frac{e^{(z+\overline\zeta)x}}{e^{2ux}+e^{2vx}}\,dx\\
 &=\frac{\pi}{2L}
     \sec\frac{\pi(z+\overline\zeta-2c)}{2L}.
\end{aligned}
\tag{MA-HY.6a}
\]
The integral converges absolutely for $z,\zeta\in\Pi$.
Write its denominator as $2e^{2cx}\cosh(Lx)$ and substitute
$y=e^{2Lx}$.  The resulting beta integral is
$\int_0^\infty y^{r-1}(1+y)^{-1}\,dy=\pi/\sin(\pi r)$,
where $0<\operatorname{Re}r<1$, proving the displayed value.
Cauchy--Schwarz against the same weight proves boundedness
of every point evaluation.  Choose the analytic square root
of $q'$ positive at zero.  The cosine addition identity gives
\[
 q'(w)^{1/2}\overline{q'(\omega)^{1/2}}\,
 K_\Pi(q(w),q(\omega))
             =\frac1{1-w\overline\omega}.
\tag{MA-HY.6b}
\]
Indeed the secant argument becomes
$\arctan w+\arctan\overline\omega$, whose cosine is
\[
 (1-w\overline\omega)/
 \sqrt{(1+w^2)(1+\overline\omega^2)}.
\]
The right side of (MA-HY.6b) is the disk $H^2$ kernel.
The weighted composition map in (MA-HY.6), with its transported
norm, therefore has exactly the disk kernel.  Equality of the
inner products on the dense spans of kernel functions proves
that its range completes to all of $H^2(\mathbb D)$.
Point evaluations identify this unitary with the stated
weighted composition map, proving both normalization and
surjectivity directly.

It commutes with multiplication in the exact sense
\[
        \mathcal U M_F\mathcal U^{-1}=M_{F\circ q}
                    \qquad(F\in H^\infty(\Pi)).
 \tag{MA-HY.7}
\]

\paragraph{MA-HY-001 (the actual closure is the interior Blaschke ideal).}
Let $B$ be the disk Blaschke product whose zeros, with their
actual multiplicities, are
\[
       \varphi(\rho-a),\qquad
       \Xi(\rho)=0,\quad u<\operatorname{Re}(\rho-a)<v.
 \tag{MA-HY.8}
\]
If there are no such zeros, take $B=1$.
Then
\[
 \mathcal U\mathcal M J=B H^2(\mathbb D),\qquad
       Z_\Sigma\simeq K_B:=H^2(\mathbb D)\ominus B H^2(\mathbb D).
 \tag{MA-HY.9}
\]
The first equality concerns the actual ideal closure in (MA-HY.1).
Equivalently it is the simultaneous vanishing, to the prescribed
multiplicities, of the Mellin jets at all zeros strictly inside
$a+\Pi$.  Zeros on the edges and outside the closed strip impose
no further conditions in this $s=0$ completion.

\begin{proof}
Put $\Psi(z)=\Xi(a+z)$.  MT-M18 (M18.6) shows that $\Psi$ is
bounded on the closed strip, and $\Psi$ is entire.  Multiplication
by $\Psi$ is therefore bounded on $\mathscr H^2(\Pi)$.
The map $T_a$ is an automorphism of $\mathcal E$, and
\[
       T_a(\psi_\eta*f)=(T_a\psi_\eta)*(T_af).
\]
The inclusion $\mathcal E\to H_\Sigma$ is continuous and has
dense range, because it contains the compact smooth functions.
Taking the closure first in $\mathcal E$ and then in $H_\Sigma$
therefore gives exactly
\[
 \mathcal M J
     =\overline{\Psi\,\mathcal M(\mathcal E)}^{\,\mathscr H^2}
     =\overline{\Psi\,\mathscr H^2(\Pi)}^{\,\mathscr H^2}.
 \tag{MA-HY.10}
\]
Thus there is no replacement of the specified theta ideal by
an independently chosen spectral kernel.

Set $F_0=\Psi\circ q\in H^\infty(\mathbb D)$.
Its canonical factorization is
\[
       F_0=B S_\mu O,\qquad
 |S_\mu(w)|=\exp\left(-\int_{\mathbb T}
              \frac{1-|w|^2}{|\zeta-w|^2}\,d\mu(\zeta)\right),
 \tag{MA-HY.11}
\]
where $\mu$ is a finite positive measure singular with respect
to arc length, and $O$ is outer.  The unimodular constants are
irrelevant.  The existence and this precise form are Beurling's
Factorization Lemma, \S7, pp.~245--246; his Theorem~I, \S9,
pp.~247--249, says that
$\overline{F_0H^2}=B S_\mu H^2$.  Here the zeros of $B$
are exactly (MA-HY.8), so their Blaschke summability is also
part of the factorization theorem.  It remains to prove
that $\mu=0$; this is a substantive step at the two ends.

Write $M=\|F_0\|_\infty>0$.  The outer-function formula gives
$|O|\le M$, whence
\[
           \log|F_0(w)|\le\log M-P\mu(w).
 \tag{MA-HY.12}
\]
On any closed arc avoiding $\{\pm i\}$ and the boundary zeros
of $F_0$, the function $F_0$ extends holomorphically across the
arc and is bounded away from zero in a neighborhood of it.
Equation (MA-HY.12) bounds $P\mu$ in the adjacent disk sector.
This forces $\mu$ to be absolutely continuous on the arc:
for every sufficiently short subarc $I$ there, the Poisson
kernel at radius $1-|I|$ over its midpoint is at least
$C/|I|$ on $I$, so $\mu(I)\le C'|I|$.  Since $\mu$ is
singular, its restriction to that arc is zero.

The finite boundary zeros of $F_0$ are isolated away from
$\{\pm i\}$, and form a countable set.  Hence $\mu$ is
supported on that countable set together with the two ends,
and must be purely atomic.  At a finite boundary zero
$\zeta$, analytic continuation gives
$|F_0(r\zeta)|\ge C(1-r)^m$ for $r$ sufficiently close to one,
where $m$ is the finite order of the zero.  An atom
$\mu(\{\zeta\})=\tau>0$ would instead give, by (MA-HY.12),
\[
       |F_0(r\zeta)|\le M
                 \exp\left(-\tau\frac{1+r}{1-r}\right),
\]
which is incompatible with the preceding lower bound.
All such atoms vanish.

Finally MT-M18 (M18.8) supplies ordinates $h_n\to+\infty$
and $h_n\to-\infty$ such that, at the fixed real part $a+c$,
\[
       \log|\Xi(a+c+ih_n)|
                 \ge -C(1+|h_n|)-C\log(2+|h_n|).
 \tag{MA-HY.13}
\]
These are unconditional selected-height minimum-modulus bounds.
By (MA-HY.5),
\[
       \varphi(c+it)=i\tanh\frac{\pi t}{2L},\qquad
       1-\left|\varphi(c+it)\right|
                    \asymp e^{-\pi|t|/L}.
\]
An atom at $i$ or at $-i$ would therefore make the upper
bound in (MA-HY.12), along the corresponding sequence,
at most $\log M-C'\exp(\pi|h_n|/L)$.  This contradicts
(MA-HY.13).  Thus $\mu=0$.  Equations (MA-HY.7),
(MA-HY.10) and Beurling's closure theorem prove (MA-HY.9).

Divisibility by $B$ in $H^2$ is equivalent to vanishing
at its zeros with their multiplicities, by the same
canonical factorization.  On a boundary line, in contrast,
the formal evaluation of a Mellin derivative of order $j$
would have squared dual norm
\[
    \int_{\mathbb R}
       \frac{|x|^{2j}e^{2ux}}{e^{2ux}+e^{2vx}}\,dx=\infty
\]
at the left edge, and the analogous integral is infinite
at the right edge.  Riesz representation, tested on compact
smooth functions, shows that none of these boundary jets
is continuous.  The closure calculation already proves that
no additional boundary restriction survives.
\end{proof}

\subsection{Compressed multipliers and the partial trace}

For $F\in H^\infty(\Pi)$ write
\[
       T_F=P_{K_B}M_{F\circ q}\big|_{K_B}.
 \tag{MA-HY.14}
\]
This is the quotient multiplier, represented on the orthogonal
model space.  In general $M_{F\circ q}$ itself does not leave
$K_B$ invariant.  Because $BH^2$ is invariant under every
bounded analytic multiplier, the compressed operators satisfy
\[
               T_{F_1}T_{F_2}=T_{F_1F_2}.
 \tag{MA-HY.15}
\]
Indeed the omitted term in inserting $1=P_{K_B}+P_{BH^2}$
has range in $BH^2$, and its final projection to $K_B$ is zero.

\paragraph{MA-HY-002 (trace class without a uniform edge margin).}
For every $g\in C_c^\infty(\mathbb R)$, the actual quotient
convolution $C_g^{Z_{\max}}$ is trace class and
\[
       \operatorname{Tr}_{Z_{\max}}C_g^{Z_{\max}}
       =\sum_{\substack{\Xi(\rho)=0\\
                  u<\operatorname{Re}(\rho-a)<v}}
                     m_\rho\,\mathcal Mg(\rho-a).
 \tag{MA-HY.16}
\]
The series is absolutely convergent.  The result also holds
for $g\in\mathcal E$.  For every real $\kappa>v$ and integer
$k\ge2$, the actual quotient translation generator $A_Z$
satisfies
\[
                    (\kappa-A_Z)^{-k}\in\mathcal S_1(Z_{\max}).
 \tag{MA-HY.17}
\]
All these assertions are unconditional for the one fixed
Hecke function.

\begin{proof}
List the zeros $\lambda_j=\rho_j-a$ strictly inside $\Pi$,
with repetitions for multiplicity, and put $b_j=\varphi(\lambda_j)$.
Let $B_n$ be the product of the first $n$ normalized Blaschke
factors.  The functions
\[
       e_j(w)=B_{j-1}(w)
               \frac{\sqrt{1-|b_j|^2}}{1-\overline b_j w}
 \tag{MA-HY.18}
\]
form an orthonormal basis of $K_B$.  To verify this without
any separation condition, use
$K_{UV}=K_U\oplus U K_V$ for inner functions $U,V$.
A single Blaschke factor has one-dimensional model space
spanned by the displayed normalized reproducing kernel.
This proves the finite-product assertion inductively.
Moreover $\bigcap_nB_nH^2=BH^2$: division by each $B_n$
preserves the boundary norm, and a weakly convergent
subsequence of these bounded quotients, or canonical
factorization, gives division by $B$ in $H^2$.
Taking orthogonal complements proves completeness in the
infinite case.  Repeated zeros cause no change in the argument.

Since $|B_{j-1}|=1$ almost everywhere on the circle,
the boundary density $|e_j|^2\,d\theta/(2\pi)$ is exactly
harmonic measure at $b_j$.  Transferring it to the strip,
for $\lambda_j=u+x_j+it_j$ with $0<x_j<L$, its two
densities on the edges are
\[
 \begin{aligned}
 P_{x_j,0}(t-t_j)&=
  \frac{\sin(\pi x_j/L)}
       {2L\{\cosh(\pi(t-t_j)/L)-\cos(\pi x_j/L)\}},\\
 P_{x_j,L}(t-t_j)&=
  \frac{\sin(\pi x_j/L)}
       {2L\{\cosh(\pi(t-t_j)/L)+\cos(\pi x_j/L)\}}.
 \end{aligned}
 \tag{MA-HY.19}
\]
These formulas follow by the disk Poisson kernel and
(MA-HY.5); their combined integral is one.  In particular,
for $|t-t_j|\ge1$ both densities are bounded by
$C_L e^{-\pi|t-t_j|/L}$, uniformly for every
$0<x_j<L$, including arbitrarily close to either edge.

Suppose $F\in H^\infty(\Pi)$ has boundary bounds
$|F(u+it)|+|F(v+it)|\le C_N(1+|t|)^{-N}$.
For $|t_j|\ge2$, divide the integral in (MA-HY.19) into
$|t-t_j|\le |t_j|/2$ and its complement.  On the first
set use $|t|\ge|t_j|/2$ and total harmonic mass at most
one.  On the complement use the exponential tail and
$\|F\|_\infty$.  Enlarging the constant handles
$|t_j|<2$ and gives
\[
 \|M_{F\circ q}e_j\|_{H^2}^2
       =\int_{\partial\Pi}|F|^2\,d\omega_{\lambda_j}
       \le C_{F,N,L}(1+|t_j|)^{-2N}.
 \tag{MA-HY.20}
\]
The constant has no inverse edge-distance factor.
MT-M18's unconditional local zero count is
$O(\log(2+|T|))$ in each fixed-height window, including
multiplicity.  Consequently (MA-HY.20) is summable
over $j$ for any $N>1/2$.  Thus
$M_{F\circ q}:K_B\to H^2$ is Hilbert--Schmidt, and so
is its compression $T_F$.

Put $G=\mathcal Mg$.  Its values on the closed strip
decay faster than every power, by integration by parts.
Choose a real $d>\max(|u|,|v|)+1$ and factor
\[
       G(z)=F_1(z)F_2(z),\qquad
       F_1(z)=\frac1{d^2-z^2},\quad
       F_2(z)=(d^2-z^2)G(z).
 \tag{MA-HY.21}
\]
Both factors are bounded analytic on $\Pi$, with
boundary decay sufficient for (MA-HY.20).
They give two Hilbert--Schmidt compressed multipliers.
Equation (MA-HY.15) then gives
$T_G=T_{F_1}T_{F_2}\in\mathcal S_1(K_B)$.
On the actual Mellin quotient convolution is
multiplication by $G$, by Fubini on the dense smooth
subspace and bounded extension.  Equations (MA-HY.2),
(MA-HY.4), (MA-HY.7) and (MA-HY.9) therefore identify
this trace-class operator with the specified quotient
convolution, up to bounded similarity.

Its trace can now be computed in the actual orthonormal
basis (MA-HY.18), since trace class has already been proved.
The analytic Poisson formula gives
\[
 \langle T_Ge_j,e_j\rangle
      =\int_{\mathbb T}(G\circ q)|e_j|^2\,\frac{d\theta}{2\pi}
      =(G\circ q)(b_j)=G(\lambda_j).
 \tag{MA-HY.22}
\]
The local zero count and the rapid vertical decay also
give absolute convergence of these diagonal values.
Summing proves (MA-HY.16), with all multiplicities.

Finally the actual translation group on $H_{\max}$ has
$\|U_t\|\le e^{vt}$ for $t\ge0$.  Its quotient has the same
bound, so for $\kappa>v$ its resolvent is the operator-norm
convergent integral
$\int_0^\infty e^{-\kappa t}U_t\,dt$.
Its Mellin multiplier is $(\kappa-z)^{-1}$.
This multiplier has the $N=1$ decay required above,
so its compression is Hilbert--Schmidt.
Its square is trace class by (MA-HY.15); higher powers
are trace class by multiplication with a bounded
resolvent.  Bounded similarity gives (MA-HY.17).
\end{proof}

\subsection{The actual ambient lift and its partial geometric limit}

\paragraph{MA-HY-003 (the geometric-order cutoff trace at $s=0$).}
Take $a=1/2$, $u=-\alpha$, $v=\alpha$, with any $\alpha>0$,
and use the actual weight $w=e^{2\alpha|x|}$.
Let $P_Z$ and $P_G=1-P_Z$ be the orthogonal projections in
$H_{\max}$ onto $J^\perp$ and $J$, respectively.
For $g\in C_c^\infty(\mathbb R)$,
$C_gP_Z$ is trace class and
\[
 \operatorname{Tr}_{H_{\max}}(C_gP_Z)
  =\sum_{|\operatorname{Re}\rho-1/2|<\alpha}
                    m_\rho\,\mathcal Mg(\rho-1/2).
 \tag{MA-HY.23}
\]
Let $\chi_R\in C_c^\infty(\mathbb R)$ have uniformly bounded
suprema and converge to one almost everywhere.
Then
\[
       \|(M_{\chi_R}-1)C_gP_Z\|_1\longrightarrow0.
 \tag{MA-HY.24}
\]
If $g(0)=0$, every $M_{\chi_R}C_gP_G$ is trace class and
\[
 \lim_R\operatorname{Tr}_{H_{\max}}(M_{\chi_R}C_gP_G)
   =-\sum_{|\operatorname{Re}\rho-1/2|<\alpha}
                    m_\rho\,\mathcal Mg(\rho-1/2).
 \tag{MA-HY.25}
\]
This is a limit of traces of the actual operators in the
specified order, without any zero-location premise.

\begin{proof}
The operator part of MA-LIFT-002 uses arithmetic only to
obtain a trace-class power of the actual quotient resolvent.
MA-HY-002 provides precisely this premise here, already
with power two.  To make the replacement explicit, let
$L=w^{-1}\partial_x(w\,\cdot)$ generate the actual adjoint
translation group, put $R=(\kappa-L)^{-1}$ for
$\kappa>\alpha+1$, and extend its restriction to
$J^\perp$ by $P_Z$.  Equation (MA-HY.17) gives
$R^2P_Z\in\mathcal S_1(H_{\max})$.
It maps boundedly into the graph domain of $L^2$.
The integration by parts on the two half-lines in
MA-LIFT.12--MA-LIFT.16, with $k=2$ and $s=0$, is the
exact identity
\[
        C_gP_Z=B_{g,2}R^2P_Z+
              \sum_{m=0}^{1}
                \phi_{g,2,m}\otimes(\ell_mR^2P_Z).
 \tag{MA-HY.26}
\]
Here $B_{g,2}$ is bounded, $\phi_{g,2,m}$ is compact
smooth, and the graph-domain trace functionals
$\ell_m(u)=\partial^m(wu)(0)$ are bounded after
$R^2P_Z$.  The finite boundary contribution at the
cusp of $w$ is retained.  Thus (MA-HY.26) proves trace
class of the actual lift.  Trace cyclicity with $P_Z$
identifies its trace with the actual quotient trace,
which MA-HY-002 computes as (MA-HY.23).

The proof of MA-LIFT-003 is purely an operator calculation
once this trace-class lift is available.
Uniform boundedness and strong convergence of the
$M_{\chi_R}$ imply (MA-HY.24) by finite-rank approximation.
The smooth compact kernel of $M_{\chi_R}C_g$ has trace
$g(0)\int\chi_R$, also for the weighted space by bounded
local conjugation and cyclicity.  Subtracting
$M_{\chi_R}C_gP_Z$ proves trace class of
$M_{\chi_R}C_gP_G$ and, for $g(0)=0$, (MA-HY.25).
No convergence of $M_{\chi_R}C_gP_G$ itself in trace norm
is asserted.
\end{proof}

\paragraph{Scope and the remaining boundary problem.}
This is an unconditional partial-strip theorem.  It does
not identify (MA-HY.25) with the full Meyer geometric trace:
the excluded zeros, including any on the two edges, are
absent from this actual $s=0$ quotient.  The difference from
the full spectral trace is exactly their absolutely
convergent Mellin sum.  MA-HM-003 prevents a full return
intertwiner for arbitrary tests if a strictly exterior
zero exists.  Under the single-character RH, every positive
centered strip contains the full divisor and the earlier
actual Meyer transport applies.

The condition $s=0$ is essential to the proof given here.
For positive polynomial weights the Mellin boundary space
has Sobolev regularity; at the thresholds established in
MA-TR-103, boundary jets are continuous and can survive in
the actual quotient.  Multiplication by the disk
half-derivative is then not the plain Hardy-space model
used above.  No assertion about the full Hardy--Sobolev
closed ideal or its trace is obtained by replacing $s=0$
by $s\ge3$ in these statements.

\paragraph{Placement in the requested geometric programme.}
The object completed here is the actual spectral jet quotient
of MT-M18/MT-M24.  The inner function $B$ encodes the divisor
that remains visible to this Hilbert norm, with all its interior
multiplicities; its Hardy-space model supplies an analytic
realization of that quotient and its trace.  This is compatible
with the nuclear trace and determinant perspective attached in
BR-TR-01/02 and F-SM-11/12 to cote~1, scans~079--080 and
111--112.  Those archival locators are inherited framework
placements, not a new independent rereading of the scans in
this section.  The Hardy model supplies neither a return to
discarded zero blocks nor an order or positivity on a
Grothendieck group.  In particular the inner factor is not an
assertion that the entire divisor is critical-line supported,
and the quotient construction does not remove the
zero-location premise from the full Meyer transport.

\paragraph{Primary literature and inherited interfaces.}
Beurling, A. (1949).
\emph{On two problems concerning linear transformations in
Hilbert space}. Acta Mathematica, 81, 239--255.
doi:10.1007/BF02395019.
The invoked locators are the Factorization Lemma in \S7,
pp.~245--246, and Theorem~I with proof in \S9,
pp.~247--249.
Bakan, A., \& Kaijser, S. (2007).
\emph{Hardy spaces for the strip}.
Journal of Mathematical Analysis and Applications, 333(1),
347--364. doi:10.1016/j.jmaa.2006.10.088.
The invoked strip-space locators are equations~(2.1)--(2.3),
Theorems~2.1--2.2 and Corollary~2.1, pp.~349--350;
the conformal/Poisson discussion is \S3 and the proofs of
the two isomorphism theorems are \S4, pp.~352--355.
The arithmetic premises are MT-M18 (M18.6), (M18.8) and
its unconditional local zero count, together with MT-M24's
identification of the actual theta ideal.
No unpublished outside result or zero-location input is
used in MA-HY-001--003.
